\renewcommand{\ElicitationDate}{2026.11.05}

\section{Elicitation \ElicitationDate}

\begin{tabular}{ll}
\textbf{Date:} & \ElicitationDate \\
\textbf{Consultant:} & Hector Hill / Gy'ax \\
\textbf{Elicitors:} & Zahra, Andrew, Matthias \\
\textbf{Location:} & TFS elicitation room \\
\textbf{Context:} & Seventh elicitation for LING 431/531 Field Methods \\
\textbf{Audio:} & \texttt{November\_05\_2026\_HH.WAV} \\
\end{tabular}

\subsection{Andrew [00:30-36:20]}

letters for reference x̱ḵg̱

[00:32]
Hector asks to do the next elicatation metadata in Gitksanimx
\begin{exe}
\ex \begin{xlist}
\item \gll
hlox̱s-a sa'nax̱\\
month-\Attr{} set.traps\\
\glt "November" [lit: the month to set traps] \hfill HH.\ElicitationDate{}
\item \gll
xwsdins sa tun ii hlox̱s-a sa'nax̱\\
five day this and month-\Attr{} set.traps\\
\glt "It's the 5th of November" \hfill HH.\ElicitationDate{}
\item \gll
Andrew=hl we-'y ii Zahra=hl we-'n ii Matthias=hl we-'n\\
Andrew=\Cn{} name-1.\Ii{} and Zahra=\Cn{} name-2.\Ii{} and Matthias=\Cn{} name-2.\Ii{}\\
\glt "My name is Andrew and Zahra is their name and Matthias is their name" \hfill HH.\ElicitationDate{}
\end{xlist}
\end{exe}

[3:17] --> re-eliciting for Zahra \textit{we're cooking a stew that needs 3 carrots, and I only brought you one. You tell me:}
\begin{exe}
\ex \gll
gilbil, win g̱al hlibuu=hl ḵawts \\
two, \Comp{} too few=\Cn{} carrots\\
\glt "Two that the carrots are too few" (SF)
\end{exe}

[3:34] --> HH: "yeah you could say that"
\begin{exe}
\ex \gll
gilbil win sg̱ooks-xw=hl ḵawts \\
two \Comp{} lack-\Val{}=\Cn{} carrots \\
\glt "Two that the carrots lack" \hfill HH.\ElicitationDate{} (VF)
\end{exe}

I'm taking this sg̱ooks translation from FV \textit{luu sḵooksxw} "to not have enough of something"; maybe the \textit{luu} is adding important information, but as it seems like the regular pre-verb, I'd be surprised

[4:03] --> seeing if Henry's hypothesis that I just needed to add a connective to make 10/29/26 [51:25] grammatical is true
\begin{exe}
\ex \gll
asag̱a-'y=hl gilbil ḵawts \\
want-1.\Ii{}=\Cn{} two carrots\\
\glt "I want two carrots" \hfill HH.\ElicitationDate{} (SF)
\end{exe}

HH: "Yes", but I realise I need to go back and try \textit{hilt} and not just \textit{ḵawts}\dots

[4:14] --> Matthias testing for extraction?
\begin{exe}
\ex \gll
* gilbil win g̱al hlibuu=hl an ḵawts?\\
{} two \Comp{} too few=\Cn{} \Ax{} carrots\\
\glt "Two that the carrots are too few" \hfill HH.\ElicitationDate{}
\end{exe}

HH: "No, just \textit{ḵawts}, \textit{en ḵawts} is where you plant the \textit{ḵawts}"

I'm assuming \textit{en} is the preposition \textit{e} combined with something else\dots so it would be "from carrots", maybe worth pursuing further?

[5:20] --> testing why Hector said people would be different in a "too few" construction (cf. 10/29/26 [45:29]) \textit{We're going to a restaurant. They have a minimum seating requirement for a table of 4 people. There’s only 3 of us. The server says:}
\begin{exe}
\ex \gll
k'i'iy win hlibuu-si'm \\
one \Comp{} few-1Pl.\Ii{} \\
\glt "We are one too few" \hfill HH.\ElicitationDate{} (VF)
\end{exe}

Very clear decomposition.

[5:45] HH: "Or you could say"
\begin{exe}
\ex \gll
Nee-dii hogax̱ gabii-si'im \\
\Neg{}-\Foc{} proper \Cnt{}-1Pl.\Ii{} \\
\glt "There's not enough of you guys" \hfill HH.\ElicitationDate{} (VF)
\end{exe}

[6:00] AS: "Can I say"
\begin{exe}
\ex \gll
moja gabii-si'im \\
almost \Cnt{}-1Pl.\Ii{}\\
\glt \textit{Intended: "there's almost enough of us"} hfill HH.\ElicitationDate{} (SF)
\end{exe}

[6:11]
\begin{exe}
\ex \gll
moja hogax̱=hl gabii-si'im\\
almost proper=\Cn{} \Cnt{}-1Pl.\Ii{}\\
\glt "There's almost enough of us" \hfill HH.\ElicitationDate{} (VF)
\end{exe}

[6:20] HH: "Our elders used to say"
\begin{exe}
\ex \gll Nee-dii am=hl moja\\
\Neg{}-\Foc{} good=\Cn{} almost\\
\glt "almost is not good" \hfill HH.\ElicitationDate{} (VF)
\end{exe} 

[6:53] --> \textit{we're splitting a pie and everyone gets a quarter}
\begin{exe}
\ex \gll
* k'i'iy hlak'aphl tx̱alpx̱ \\
{} one part.of four \\
\glt \textit{Intended: "one part of four"}  \hfill HH.\ElicitationDate{} (SF)
\end{exe}

HH: "So what are you saying on either side of this?"

[7:14]
\begin{exe}
\ex \begin{xlist}
\item \gll k'i'iy=hl \\
one \\
HH: "Now you have one of four"
\item \gll
gubi-'y=hl ḵam ki'i'yhl \\ 
eat-1.\Ii{}=\Cn{} only one=\Cn{} \\
HH: "I ate only one piece?"
\item \gll 
gubi-'y k'i'iy hlak'aphl tx̱alpx̱\\
eat-1.\Ii{} one part.of four\\
AS: \textit{Intended: "I ate one part of four"} (SF)
\item  \gll 
gubi-'y k'i'iy=hl tx̱alpx̱=hl anaax\\
eat-1.\Ii{} one=\Cn{} four=\Cn{} bread \\
"I ate one fourth of the bread" (VF)
\end{xlist}
\end{exe}

Hector confirmed that this would be only one loaf that is split into four pieces, and not eating an entire loaf when there are four loaves of bread. However, he didn't insert of \textit{hlak'aphl}, but he did accept it in the next sentence. I think the previous omission was just a slip-up.

[8:45] --> Hector says \textit{eksdam anaax} would be a bit better
\begin{exe}
\ex \gll 
gubi-'y k'i'iy=hl hlak'aphl tx̱alpx̱=hl eksda-m anaax\\
eat-1.\Ii{} one=\Cn{} part.of four=Cn{} sweet-\Attr{} bread\\
\glt "I ate one part of four of the cake" \hfill HH.\ElicitationDate{}
\end{exe}

[9:40]
\begin{exe}
\ex \gll
gubi-'y k'i'iy=hl e=hl gwile'l=hl eksda-m anaax \\
eat-1.\Ii{} one=\Cn{} from=\Cn{} three=\Cn{} sweet-\Attr{} bread \\
\glt "I ate one part of three of the cake" \hfill HH.\ElicitationDate{} (VF)
\end{exe}

[9:55] --> the following volunteered form was produced with much more hesitation than the form above
\begin{exe}
\ex \gll 
k'i'iy=hl hlak'aphl gwile'l=hl goja am eksda-m anaax \\
one=\Cn{} part.of three=\Cn{} ?? good sweet-\Attr{} bread \\
\glt "One part of three of the cake" \hfill HH.\ElicitationDate{}
\end{exe}

HH: "\textit{goja} means you cut in 3" <-- I've never seen this anywhere else and I'd like a second source on it\dots

[10:46] --> Hector trails off here as if unsure how to continue before asking for my \textit{hlak'aphl} version
\begin{exe}
\ex \gll 
gubi-'y k'i'iy=hl goja e=hl\dots \\
eat-1.\Ii{} one=\Cn{} ?? \Prep=\Cn{}\\
\end{exe}

[11:37]
\begin{exe}
\ex \gll 
gubi-'y k'i'iy=hl hlak'aphl gwile'l=hl eksda-m anaax\\
eat-1.\Ii{} one=\Cn{} part.of three=\Cn{} sweet-\Attr{} bread\\
\glt "I ate one third of the cake" \hfill HH.\ElicitationDate{}
\end{exe}

[12:05] --> Hector mentions how he arrived early this time (~11:30)
\begin{exe}
\ex \gll 
ii heksamox guks ba'a gaa=hl wilb'y\\
and again back drive? ??=\Cn{} house-1.\Ii{}\\
\glt "I drove back to my house" \hfill HH.\ElicitationDate{}
\end{exe}

\textit{baahan} is the FV entry for drive, and I don't hear any <h>, so it could be a dialectical thing?

[12:50] --> trying out 3/4
\begin{exe}
\ex \gll 
gub-'m gwile'l=hl hlak'aphl tx̱alpx̱=hl eksda-m anaax \\
eat-1Pl.\Ii{} three=\Cn{} part.of four=\Cn{} sweet-\Attr{} bread\\
\glt "We ate three quarters of the cake" \hfill HH.\ElicitationDate{} (accepted suggested form)
\end{exe}

In talking about chance/odds/possibility, I mention \textit{aḵ} as a Nisgha word for "impossible", and Hector mentions that it could also mean "open"
[14:30]
\begin{exe}
\ex \gll 
ḵ'aḵ aa'ts'ip\\
open door\\
\glt "Open the door!" \hfill HH.\ElicitationDate{}
\end{exe}

so \textit{ḵ'ak} and \textit{aḵ} are not the same words, but I can see the confusion\dots I must have had very glottal vowel starts that day

The point still stands that \textit{aḵ} is attested on FV and maybe Hector just doesn't use this word, maybe worth trying with Vince

Moving over to chance as a type of fraction, we end up on the context of a coin flip

\textit{Context: We're gambling on a coin flip, heads: I win, tails: you win. Zahra asks what the odds of winning are.}

[16:20]
\begin{exe}
\ex \gll
dim ga'a-'m ji-da x̱staa-n \\
\Fut{} see-1Pl.\Ii{} \Irr{}-\Spt{} win-2.\Ii{}\\
\glt "We'll see if you win" \hfill HH.\ElicitationDate{}
\end{exe}

[16:50] --> AS: "Can I say something like"
\begin{exe}
\ex \gll
k'i'iy luu gilbil dim x̱staa-n \\
one in two \Fut{} win-2.\Ii{}\\
\textit{Intended: "One in two you will win"} \hfill HH.\ElicitationDate{}
\end{exe}

HH: "That doesn't really say that you have a chance. "

[17:40] Trying out a new context where chance can make sense: \textit{Context: How would you tell somebody it might rain if you see clouds on the horizon?}
\begin{exe}
\ex \gll 
dim wis-ima'a \\
\Fut{} rain-\Epis{} \\
\glt "It might rain" \hfill HH.\ElicitationDate{}
\end{exe} 

HH: "There's got to be a word between \textit{k'i'iy} and \textit{x̱staa}"

[19:00]
\begin{exe}
\ex begin{xlist}
\item \gll 
yugw-imaa \\
\Prog{}-\Epis{} \\
\glt "Maybe"
\item \gll
dim beligii x̱staa' 'niin \\
\Fut{} suddenly win 2.\Iii{} \\
\glt "You might win" \hfill HH.\ElicitationDate{}
\end{exe}

Looking back at this after finishing, \textit{beligi} could be potentially decomposed into \textit{be} and \textit{ligi}, but that would mean explain what kind of modifier the former morpheme is, and it doesn't look to be too productive.

[19:25] --> AS: "Could I say"
\begin{exe}
\ex \gll 
an-yugw-imaa \\
\Nmlz{}-\Prog{}-\Epis{}\\
\glt \textit{Intended: a nominalized version of maybe to mean "chance"}
\end{exe}

[19:44] HH: "If I give you the winning numbers in a envelope then"
\begin{exe}
\ex \gll 
yugw-imaa x̱staa-n\\
\Prog{}-\Epis{} win-2.\Ii{}\\
\glt "You might win" \hfill HH.\ElicitationDate{}
\end{exe}

[20:33] --> ZA: "What about"
\begin{exe}
\ex \gll
k'i'iy luu gilbil win ji-x̱staa-'y \\
one in two \Comp{} \Irr{}-win-1.\Ii{}\\
\glt \textit{Intended: "one in two that I might win"} \hfill HH.\ElicitationDate{}
\end{exe}

[22:50] --> two options that were come up with
\begin{exe}
    \ex \begin{xlist}
    \item \gll 
    k'i'iy ksi gilbil ji-da x̱staa' \\
    one out.of two \Irr{}-\Spt{} win \\
    \glt \textit{Intended: "one out of two you might win"}
    \item \gll 
    k'i'iy e=hl gilbil ji-da x̱staa'\\
    one \Prep=\Cn{} two \Irr{}-\Spt{} win \\
    \glt \textit{Intended: "one from two you might win"}
    \end{xlist}
\end{exe}

HH: "We need to find the word for flip, because that's an important part of whether you win or lose"

[29:48]
\begin{exe}
\ex \gll 
k'i'iy win bax-hileldin=hl dildilaax̱xw-m gwiikw \\
one \Comp{} up-flip=\Cn{} silver-\Attr{} marmot \\
\glt "one flip of a coin" \hfill HH.\ElicitationDate{}
\end{exe}

[30:17] HH: "\textit{dimin gyaahl jida x̱staa} at the end. Now this is throwing up the coin and you only have one chance to win" 

[31:55]
\begin{exe}
\ex \gll 
(ḵ'am) k'i'iy win bax-hileldin=hl dildilaax̱xw-m gwiikw dim-im gaa=hl ji-da x̱staa-n \\
(only) one \Comp{} up-flip=\Cn{} silver-\Attr{} marmot \Fut{}-2.\I{} see=\Cn{} \Irr{}-\Spt{} win-2.\Ii{}\\
\glt "If you flip a coin the chance you have of winning is one" \hfill HH.\ElicitationDate{}
\end{exe}

This construction took 6 minutes to put together and it's not even using odds!! It's fine\dots

[33:04] AS: "How would I say the odds are one in two of winning?"
\begin{exe}
\ex \gll
(ḵ'am) k'i'iy ji-ligi gilbil \\
(only) one \Irr{}-\Dwid{} two \\
\glt "one or two" \hfill HH.\ElicitationDate{} (VF) (VG) 
\end{exe}

ZA offers or as a translation of \textit{ligi}, but I'm not sure where this comes from? Hector agress, but I'm glossing it as a domain widener a la the grammatical sketch; further investigation is warranted. \\

--> Yeah I don't think "or" would be a good translation, I just thought he might associate it with or bc it very often gets used with the actual word for or \textit{oo}, as \textit{oo ligi} -- from BDM. Slightly more conventional translation of \textit{ligi} as "any" might work too I suppose -- "one of any two"

[33:56] --> Trying a dice context. \textit{Context: we're rolling a die and you only win if you land on a 6. The odds are:}
\begin{exe}
\ex \gll 
(ḵ'am) k'i'iy ji-ligi ḵ'oo'lt\\
(only) one \Irr{}-\Dwid{} six \\
\glt "One out of six"
\end{exe}

HH: "Yep, that works"

[35:30] Hector's construction for casino, for fun
\begin{exe}
\ex \begin{xlist}
\item \gll 
wilp-'m xstaa imaa\\
house-\Attr{} win \Epis{}\\
\item \gll 
wilp win jehl-m \\
house \Comp{} lose-\Attr{}\\
\item \gll 
wilp-'m moja xstaa \\
house-\Attr{} almost win \\
\end{xlist}
\end{exe}

\subsection{Zahra [36:20-1:07:40]}

[36:53]
\begin{exe}
    \ex \glll
    sdemḵhl ben’y \\
    sdemḵ=hl ben-’y \\
    make.noise=\Cn{} belly-\Fsg.\Ii \\
    \glt "My stomach is growling, making noise." \hfill HH.\ElicitationDate{} (VF)
\end{exe}

[37:12]
\begin{exe}
    \ex \glll
    x̱dex ’nii’y \\
    x̱dax ’nii’y \\
    hunger \Fsg.\Iii \\
    \glt "I'm hungry." \hfill HH.\ElicitationDate{} (VF)
\end{exe}

[37:14]
\begin{exe}
    \ex \glll
    x̱a’nii ’niin x̱deyi’yaa? \\
    x̱’ni ’niin x̱dex-’y=aa\\
    hear \Ssg.\Iii{} hungry-\Fsg.\Ii=\Q \\
    \glt "Did you hear that I'm hungry?" \hfill HH.\ElicitationDate{} (VF)
\end{exe}

\textit{xʸ} \(\rightarrow\) \textit{y} before a sonorant suffix is shown with this root in the stories, e.g. \textit{x̱deyi'm} with -\Fpl.\Ii. Not totally sure about the question particle, it sounded a bit like \textit{e} rather than \textit{aa}, could be worth checking. 

[39:40]
\begin{exe}
    \ex\glll
    (g̱ay) k'aa *(hilthl) haanak' ehl get begwit go'ohl luu gwendinst \\
    (g̱ay) k'aa hilt=hl haanak' a[-t]=hl gat bakw-it go'o=hl luu gwendinst \\
    \Cntr{} more many=\Cn{} woman \Prep{}[-3.\Ii]=\Cn{} man come-\Sx{} \Loc=\Cn{} in event \\
    \glt "More women than men came to the party." \hfill HH.\ElicitationDate{} (VF) \\
    Comment: "Gotta have the number, or \textit{hilt}." \\
    Comment: Preferred with \textit{g̱ay} but still grammatical without. 
\end{exe}

HH comments about \textit{g̱ay} being preferred: "you're literally separating out that there are more women." Also seems to say that it's grammatical to drop the PP referring to the men (i.e. just, "more women came to the party") as the comparison with the number of men is clear with the context.

[45:40]
\begin{exe}
    \ex \glll
    ’Wii ’nakwt Andrew (ii hlep dulpxwin). \\
    ’wii ’nakw-t Andrew ii hlep dulp-xwin. \\ 
    big long=\Pn{} Andrew \Ccnj{} deep short-\Val-\Ssg.\Ii{} \\
    \glt "Andrew is tall (and you're short)." \hfill HH.\ElicitationDate{} (SF)
\end{exe}

[46:29]
\begin{exe}
    \ex \glll
    enyentxw'm g̱oot Matthias \\
    enyentxw'm g̱oot Matthias \\
    funny-\Fpl.\Ii{} heart Matthias \\
    \glt "Matthias is funny." \hfill HH.\ElicitationDate{} (VF)
\end{exe}

Above is potentially something like "Matthias is funny to us," etc. Otherwise possible that there's -\textit{m} \Attr{} but then some facts about predicative status might be unclear. 

[47:10]
\begin{exe}
    \ex \glll
    guxwhl luu amhl g̱oots \\
    guxhl luu amhl g̱oots \\
    \Hab=\Cn{} in good=\Cn{} heart[-3.\Ii]=s \\
    \glt "They're always happy." \hfill HH.\ElicitationDate{} (VF)
\end{exe}

[47:23]
\begin{exe}
    \ex \glll
    dim duu’w amhl goot get ii dest gat hlag̱ats’uut \\
    dim duu'w am=hl goot gat ii des=t gat hla-g̱a-ts’uu-t \\
    \Prosp{} far good=\Cn{} heart person \Ccnj{} touch=3.\I person \Part-\Distr-other-3.\Ii \\
    \glt "When you have a good heart, or a funny heart... it touches other people." \hfill HH.\ElicitationDate{} (VF)
\end{exe}

Lots I'm unsure about here, especially the beginning before \textit{amhl} and whatever is between \textit{des} and \textit{hlag̱ats’uut} (which could instead be e.g. something preposition-y).

[47:53]
\begin{exe}
    \ex \glll
    get g̱adaaxt \\
    get g̱a-taax-t \\
    person \Distr-perimeter-3.\Ii \\
    \glt "surrounding people" \hfill HH.\ElicitationDate{} (VF)
\end{exe}

[49:55]
\begin{exe}
    \ex \glll
    moomst Matthias \\
    moomst[-t?] Matthias \\
    crazy[-3.\Ii] Matthias \\
    \glt "Matthias is crazy, silly." \hfill HH.\ElicitationDate{} (SF)
\end{exe}

[51:33]
\begin{exe}
    \ex \glll
    ax̱g̱alḵ'an 'niin \\
    ax̱g̱alḵ'an 'niin \\
    immature \Ssg.\Iii \\
    \glt "You're immature." \hfill HH.\ElicitationDate{} (VF)
\end{exe}

FV sites differ on \textit{ax̱ g̱alḵ'an} `thoughtless, disobedient, immature, ignorant' being two words or one. 

[51:40]
\begin{exe}
    \ex \glll
    dax̱gedinhl g̱oodin \\
    dax̱gat-in=hl g̱ood-n \\    
    strong-\Caus2=\Cn{} heart-\Ssg.\Ii \\
    \glt "Strengthen your heart." \hfill HH.\ElicitationDate{} (VF)
\end{exe}

[52:03]
\begin{exe}
    \ex \glll
    wilaat'ikst Matthias \\
    wilaat'iks=t Matthias \\
    learn=\Pn{} Matthias \\
    \glt "Matthias is smart." \hfill HH.\ElicitationDate{} (VF)
\end{exe}

[52:30]
\begin{exe}
    \ex \glll
    mumḵ’an! \\
    mumḵ’-n\\
    smile-\Ssg.\Ii \\
    \glt "Smile!" \hfill HH.\ElicitationDate{} (VF)
\end{exe}

[52:54]
\begin{exe}
    \ex \glll
    ha'w ji mumḵ’an \\
    ha'w ji mumḵ’-n  \\
    don't \Irr{} smile-\Ssg.\Ii \\
    \glt "Don't smile" \hfill HH.\ElicitationDate{} (VF)
\end{exe}

[53:10]
\begin{exe}
    \ex \glll
    'wii nakwt Andrew ii 'malut \\
    'wii nakw=t Andrew ii 'malu=t \\
    big long=\Pn{} Andrew \Ccnj{} crazy=3.\Ii \\
    \glt "Andrew is tall and crazy." \hfill HH.\ElicitationDate{} (VF)
\end{exe}

Again should figure out whether -\textit{t} is \Sx{} with movement or \Pn{} with elision or simply -3.\Ii.

[54:28]
\begin{exe}
    \ex \glll
    Sag̱ayt k’i’yhl g̱a’nakwt Andrew ehl *(win) g̱a’malut. \\
    Sag̱ayt k’i’y=hl g̱a-’nakw=t Andrew e[-t]=hl win g̱a-’malu-t. \\
    together one=\Cn{} \Distr-long=\Pn{} Andrew \Prep{}[3.-\Ii]=\Cn{} \Comp{} \Distr-crazy-3.\Ii\\
    \glt "Andrew is as tall as he is crazy." \hfill HH.\ElicitationDate{} (SF)
\end{exe}

Metalinguistic equative. VF lacked \textit{win} but HH corrected the VF to add it.

[54:28]
\begin{exe}
    \ex \glll
    Sag̱ayt k’i’yhl g̱a’malut Matthias ehl win g̱a’malut Andrew. \\
    Sag̱ayt k’i’y=hl g̱a-’malu=t Matthias e[-t]=hl win g̱a-’malu=t Andrew. \\
    together one=\Cn{} \Distr-crazy=\Pn{} Matthias =\Prep{}[3.-\Ii]=\Cn{} \Comp{} \Distr-crazy=\Pn{} Andrew \\
    \glt \textit{Intended}: "Andrew is as tall as he is crazy." \hfill HH.\ElicitationDate{} (SF)
\end{exe}

[56:20]
\begin{exe}
    \ex \glll
    Sag̱ayt k’i’y win ’mag̱alus dip. \\
    sag̱ayt k’i’y=hl ’ma<g̱a>lu-t=s dip \\
    together one \Comp{} <Pl>crazy[-3.\Ii]=\Pn{} \Assoc \\
    \glt "They are as crazy as each other." \hfill HH.\ElicitationDate{} (VF)
\end{exe}

[56:55]
\begin{exe}
    \ex \glll
    Sag̱ayt k’i’yhl g̱a’malus dip Andrew g̱an Matthias. \\
    Sag̱ayt k’i’y=hl g̱a-’malu[-t]=s dip Andrew g̱an Matthias. \\
    together one=\Cn{} \Distr-crazy=\Pn{} \Assoc{} Andrew \Pcnj{} Matthias \\
    \glt "Andrew is as crazy as Matthias; they are as crazy as each other, Andrew and Matthias." \hfill HH.\ElicitationDate{} (SF)
\end{exe}

I didn't realize the plural infixation thing with \textit{’mag̱alus} so I kept using \textit{g̱a’malus}; not sure if HH thought I was making a speech error but he accepted this. Would make some sense for either to be acceptable in a context with \textit{dip} maybe. Would be interesting to see whether we can have \textit{g̱a’mag̱alus} with the degree nominalizer, to see how much structure (e.g. the plural) the prefix can take in its degree nominalizing use.

[57:52]
\begin{exe}
    \ex \glll
    Sag̱ayt k’i’yhl g̱asg̱oos Andrew ehl win malut. \\
    sag̱ayt k’i’y=hl g̱asg̱oo[-t]=s Andrew a[-t]=hl win malu-t. \\
    together one=\Cn{} \Mass[-3.\Ii]=\Pn{} Andrew \Prep{}[-3.\Ii]=\Cn{} \Comp{} crazy-3.\Ii \\
    \glt "Andrew is as tall as he is crazy." \hfill HH.\ElicitationDate{} (VF)
\end{exe}

[58:20]
\begin{exe}
    \ex \glll
    Sag̱ayt k’i’y win malu=s Andrew ehl g̱asg̱oot. \\
    sag̱ayt k’i’y win malu=s Andrew a[-t]=hl g̱asg̱oo-t. \\
    together one \Comp{}   crazy=\Pn{} Andrew \Prep{}[-3.\Ii]=\Cn{} \Mass[-3.\Ii]-3.\Ii \\    
    \glt "Andrew is as crazy as he is tall." \hfill HH.\ElicitationDate{} (VF)
\end{exe}

HH went through a few different versions out loud here. 

\begin{exe}
    \ex \glll
    Sag̱ayt k’i’yhl g̱asg̱oos Andrew ehl ’malut. \\
    Sag̱ayt k’i’y=hl g̱asg̱oo[-t]=s Andrew a[-t]=hl ’malu-t. \\
    together one=\Cn{} \Mass[-3.\Ii]=\Pn{} Andrew \Prep{}[-3.\Ii]=\Cn{} crazy-3.\Ii \\
    \glt "Andrew is as crazy as he is tall; Andrew's height and craziness are similar." \hfill HH.\ElicitationDate{} (VF)
    % Comment: "Now we're talking about, the height and the craziness are similar."
\end{exe}

[59:55]
\begin{exe}
    \ex \glll
    ’Wii ’nakwt Andrew ii ḵ’ap ’malut. \\
    ’wii ’nakw=t Andrew ii ḵ’ap ’malu-t. \\
    big long=\Pn{} Andrew \Ccnj{} \Verum{} crazy-3.\Ii \\
    \glt "Andrew is tall but he's crazy." \hfill HH.\ElicitationDate{} (VF)
\end{exe}

[1:01:19]
\begin{exe}
    \ex \glll
    Sag̱ayt k’i’yhl g̱asg̱oon ehl ’malun. \\
    Sag̱ayt k’i’y=hl g̱asg̱oo-n a[-t]=hl ’malu-n. \\
    together one=\Cn{} \Mass-\Ssg.\Ii{} \Prep{}[-3.\Ii]=\Cn{} crazy-\Ssg.\Ii \\
    \glt "You're as crazy as you are tall; you're the same height as your craziness." \hfill HH.\ElicitationDate{} (VF)
    % Comment: "Now we're talking about, the height and the craziness are similar."
\end{exe}

HH comments on the way of phrasing these as comparisons: "But this terminology is the Western society's."

[1:02:02]
\begin{exe}
    \ex \glll
    g̱a’negwit agwi ji ḵ’ap ’malunis  \\
    g̱a’-nakw-t agwi ji ḵ’ap ’malu-n=is  \\
    \Distr-long-3.\Ii{} what \Irr{} \Verum{} crazy-\Ssg.\Ii=\Qudd \\
    \glt ? \hfill HH.\ElicitationDate{} (VF)
\end{exe}

Wasn't sure about this, especially the parts in the middle, nor the translation.  Sounded like there wa s some glottalization in the second half of \textit{’malu} but hard to tell.

[1:02:05]
\begin{exe}
    \ex \glll
    ’wii ’nakw ’niin l ii a ḵ’ap tḵ’alyeehl 'malun \\
    ’wii ’nakw ’niin l? ii a? ḵ’ap tḵ’al-yeehl? 'malu-n \\
    big long \Ssg.\Iii{} ? \Ccnj{} \Prep? \Verum{} against-go=\Cn{} crazy-\Ssg.\Ii \\
    \glt "You're tall, but you're still crazy." \hfill HH.\ElicitationDate{} (VF)
\end{exe}

Lots unclear here! If \textit{tḵ’alyeehl} is right, that apparently expresses an unexpected increase per the stories. No idea what the apparently random lateral would be but I may have heard it in a couple of these examples. Not sure about \textit{a}/vowel sounding thing before the verum either.

[1:02:13]
\begin{exe}
    \ex \glll
    hlep dulpxwin ii a ḵ’ap malun \\
    hlep dulp-xw-n ii a? ḵ’ap malu-n \\
    deep short-\Val-\Ssg.\Ii{} \Ccnj{} \Prep? \Verum{} \\
    \glt "You're short and you're still crazy." \hfill HH.\ElicitationDate{} (VF)
\end{exe}

[1:03:30]
\begin{exe}
    \ex \glll
    ama g̱oot ’nidis ii a ḵ’ap ’malut \\
    am-a g̱oot ’nit=is ii a? ḵ’ap ’malu-t \\
    good-\Attr{} heart 3.\Iii=\Qudd{} \Ccnj{} \Prep? \Verum{} crazy-3.\Ii{} \\
    \glt "He has a good heart but he's crazy" \hfill HH.\ElicitationDate{} (VF)
\end{exe}

\textit{ama g̱oot}; cf. FV \textit{amagat} `kind', stories \textit{amaga'} `good-\Attr-see'.

[1:04:06]
\begin{exe}
    \ex \glll
    g̱ay k’aa ’malut ’nidis,  maxhlaamag̱at \\
    g̱ay k’aa ’malu-t ’nit=is,  ma=xhlaa-mag̱-a-t \\
    \Cntr{} more crazy-3.\Ii{} 3.\Iii=\Qudd{} 2.\I=alone-put-\T-3.\Ii \\
    \glt "He's (even) crazier [than he is caring], leave him!" \hfill HH.\ElicitationDate{} (VF)
\end{exe}

Sounded like \textit{ma}= was doubled in the above, not sure if that means something. 

[1:04:18]
\begin{exe}
    \ex \glll
    g̱ay k’aa ’malut ’nidis,  hets'im ligi s xhlaamag̱at \\
    g̱ay k’aa ’malu-t ’nit=is  hets'im ligi s? xhlaa-mag̱-a-t \\
    \Cntr{} more crazy-3.\Ii{} 3.\Iii=\Qudd{} just \Dwid{} ? alone-put-\T-3.\Ii \\
    \glt "He's (even) crazier [than he is tall], just leave him!" \hfill HH.\ElicitationDate{} (VF)
\end{exe}

Again some bits I'm definitely unsure of here. Not sure if I really heard this \textit{s} or it's just the \textit{xhl} of the next word.

[1:04:26]
\begin{exe}
    \ex \glll
    hets'im ligi xhlaamag̱at, ’maluyis! \\
    hets'im ligi xhlaa-mag̱-a-t ’malu[-3.\Ii]=is \\
    just \Dwid{} alone-put-\T-3.\Ii{} crazy[-3.\Ii]=\Qudd \\
    \glt "Just leave him, he's crazy!" \hfill HH.\ElicitationDate{} (VF)
\end{exe}

[1:06:00]
\begin{exe}
    \ex \glll
    ’malu ’nidis, ii ḵ'ap ama g̱oots t’is… neksxwhl! \\
    ’malu ’nit=is, ii ḵ'ap am-a g̱oot[-t]=s? t’is  neks-xw=hl? \\
    crazy 3.\Iii=\Qudd{} \Ccnj{} \Verum{} good-\Attr{} heart[-3.\Ii]=\Pn? big spouse-\Val=\Cn? \\
    \glt "He's crazy, but he has a good heart… marry him!" \hfill HH.\ElicitationDate{} (VF)
\end{exe}

Not sure about \textit{t’is}; maybe could also be \textit{dis} `time,' or some stacking of suffixes on \textit{g̱oot} as somehow \textit{g̱ootsdis}. HH also gave a version with \textit{am=hl} (with \Cn) instead of \textit{am-a} (with \Attr).

[1:06:34]
\begin{exe}
    \ex \glll
    hets'im ligi neksxw \\
    hets'im ligi neks-xw \\
    just \Dwid{} spouse-\Val{} \\
    \glt "Just marry him!" \hfill HH.\ElicitationDate{} (VF)
\end{exe}

[1:06:40]
\begin{exe}
    \ex \glll
    ama get ’nidis(, neksxw!)  \\
    am-a gat ’nit=is neks-xw \\
    good-\Attr{} man 3.\Iii=\Qudd{} spouse-\Val{} \\
    \glt "He's a good man(, marry him!)." \hfill HH.\ElicitationDate{} (VF)
\end{exe}

Hard to hear whether a \Cn{} is present at the end.

[1:06:56]
\begin{exe}
    \ex \glll
    mooji ’maludis, ii ḵ’ap amhl g̱ootsdis, ii ama getsdis, neksxwhl! \\
    mooji ’malu-t=is, ii ḵ’ap am=hl g̱oot=s=d?=is, ii am-a get=s=d?=is neks-xw=hl \\
    almost crazy-3.\Ii=\Qudd{} \Ccnj{} \Verum{} good=\Cn{} heart[-3.\Ii]=?=?=\Qudd \Ccnj{} good-\Attr{} man=?=?=\Qudd{} spouse-\Val=\Cn? \\
    \glt "He's almost crazy, but he has a good heart and he's a good man, marry him!" \hfill HH.\ElicitationDate{} (VF)
\end{exe}

\subsection{Matthias [1:07:40-1:36:00]}

[1:09:02]: Both are good, but HH says the difference is that \textit{hilt smex ga'a'y} could be anywhere, whereas \textit{hilthl smex ga'a'y} is in a certain place. For example, ``I seen a lot of bears at the fishing ground''.



\begin{exe}
    \ex
    \begin{xlist}
        \item \glll
        hilt smex ga'a'y \\
        hilt smex ga'a-'y \\
        many bear see-\Fsg{}.\Ii{} \\
        \glt "I saw many bears." \hfill HH.\ElicitationDate{} (SF)
        \item \glll
        hilthl smex ga'a'y \\
        hilt=hl smex ga'a-'y \\
        many=\Cn{} bear see-\Fsg{}.\Ii{} \\
        \glt "I saw many bears." \hfill HH.\ElicitationDate{} (SF)
    \end{xlist}
\end{exe}

[1:10:00]: Okay, so \textit{g̱a'a} is \Loc{} as per the grammatical sketch, but why is there a FV entry with \textit{asun} (\url{https://www.firstvoices.com/gitsenimx/words/4f1e937d-5a83-4dda-ab65-ef87a881d7a0}). Could \textit{g̱a'a} actually decompose differently?

Also, rigsby has \textit{an/en} as its own prefix literally meaning ``place where some activity takes place'', but honestly this is probably the nominalizer as per the grammtical sketch. I think this makes sense, since \textit{aat} is the verb for net fishing?

\begin{exe}
    \ex
    \begin{xlist}
        \item \glll
        hilt smexhl ga'a'y g̱a'asun \\
        hilt smexhl ga'a-'y g̱a'a[-t]-sun\\
        many bear=\Cn{} see-\Fsg{}.\Ii{} \Loc{}[-3.\Ii{}]=\Pn{}=\Demprox{}\\
        \glt "I saw many bears over here." \hfill HH.\ElicitationDate{} (VF)
        \item \glll
        hilt smex ga'a'y g̱a'ahl enaat \\
        hilt smex ga'a-'y g̱a'a=hl en-aat  \\
        many bear see-\Fsg{}.\Ii{} \Loc{}=\Cn{} \Nmlz{}-net.fish \\
        \glt "I saw many bears at the (net) fishing grounds." \hfill HH.\ElicitationDate{} (VF)
    \end{xlist}
\end{exe}

[1:10:40]: Again, it is rejected when I attempt with attributive \textit{-a}.

HH says its a question, and I ask if it is similar to \textit{hilt smex ga'anaa}, but idk. \textit{hilda} on its own gan be a question meaning ``lots?''

\begin{exe}
    \ex \glll
    * hilda smex ga'a'y \\
    {} hilt-a smex ga'a-'y\\
    {} many-\Attr{} bear see-\Fsg{}.\Ii{} \\
    \glt Intended: "I saw many bears." \hfill HH.\ElicitationDate{} (SF)
\end{exe}

[1:11:50]: Very interesting, in this case it is possible to use \textit{hilda} with \textit{smex}, but only with the `meat' meaning. Does this mean \textit{hilda} with attributive marking necessarily requires a mass noun?

\begin{exe}
    \ex \glll
    hilda smexhl ga'a'y \\
    hilt-a smex=hl ga'a-'y\\
    many-\Attr{} meat=\Cn{} see-\Fsg{}.\Ii{} \\
    \glt \cmark  I saw a lot of meat. \xmark  I saw many bears.  \hfill HH.\ElicitationDate{} (SF)
\end{exe}

[1:12:48]: So apparently it's still kind of broken when you put \textit{hilda} first. I guess this makes sense since if the entire phrase \textit{hilda smex} is fronted it is likely focused.

\begin{exe}
    \ex \glll
    hilda smexhl gupdiit \\
    hilt-a smex=hl gup-diit \\
    many-\Attr{} meat=\Cn{} eat-\Tpl{}.\Ii{} \\
    \glt They ate a lot of meat.  \hfill HH.\ElicitationDate{} (VF)
\end{exe}

[1:13:08]: This is a ``little bit better''. So not entirely good?

Zahra also tried kind of focusing it with the fronted \textit{hilda smexhl}, but still not fully accepted.

\begin{exe}
    \ex \glll
    gupdiit hilda smex  \\
    gup-diit hild-a smex \\
    eat-\Tpl{}.\Ii{} many-\Attr{} meat \\
    \glt They ate a lot of meat.  \hfill HH.\ElicitationDate{} (SF)
\end{exe}

[1:14:43]: This also works. I wasn't careful about elucidating the different between this and \textit{hilthl smex} above though.

\begin{exe}
    \ex \glll
    hilt smexhl ga'a'y  \\
    hilt smex=hl ga'a-'y \\
    many bear=\Cn{} see-\Fsg{}.\Ii{}  \\
    \glt I saw a lot of bears.  \hfill HH.\ElicitationDate{} (SF)
\end{exe}

[1:14:56]: HH: The question here would be ``where?''. 

Hector also gives a context for the meat reading, e.g. you are at a feast.

\begin{exe}
    \ex \glll
    hilthl smexhl ga'a'y  \\
    hilt=hl smex=hl ga'a-'y \\
    many=\Cn{} bear=\Cn{} see-\Fsg{}.\Ii{}  \\
    \glt \cmark I saw a lot of bears. \cmark I saw a lot of meat.  \hfill HH.\ElicitationDate{} (SF)
\end{exe}

[1:16:20]: Can drop \textit{=hl} here. But interpretations may not be idential; see below.

\begin{exe}
    \ex \glll
    hilt(hl) usxw'y \\
    hilt(=hl) us-xw-'y \\
    many(=\Cn{}) dog-\Val{}-\Fsg{}.\Ii{} \\
    \glt "I have many dogs." \hfill HH.\ElicitationDate{} (SF)
\end{exe}

[1:16:55]: So this can be used as a kind of assertion (using \Qudd{}). But once again HH here says that with vs. without \textit{=hl} implies a ``different place''. This should really be looked into.

\begin{exe}
    \ex \glll
    hilt usxw'yis \\
    hilt us-xw-'y=is \\
    many dog-\Val{}-\Fsg{}.\Ii{}=\Qudd{} \\
    \glt "I have many dogs." \hfill HH.\ElicitationDate{} (VF)
\end{exe}

[1:18:08]: So here when I try to drop the \textit{=hl}, HH responds with \Qudd{} again. So it looks like it's possible, not sure if entirely the same as above.

\begin{exe}
    \ex \glll
    gilbilhl usxw'y \\
    gilbil=hl us-xw-'y \\
    two=\Cn{} dog-\Val{}-\Fsg{}.\Ii{} \\
    \glt "I have two dogs." \hfill HH.\ElicitationDate{} (SF)
\end{exe}

[1:19:23]: This could be a ``conversation starter'', prompted by Andrew. 

Can respond \textit{ḵ'am gilbilaa} `only two'?

\begin{exe}
    \ex \glll
    ḵ'am gilbilhl usxw'y \\
    ḵ'am  gilbil=hl us-xw-'y \\
    onyl two=\Cn{} dog-\Val{}-\Fsg{}.\Ii{} \\
    \glt "I only have two dogs." \hfill HH.\ElicitationDate{} (VF)
\end{exe}

[1:19:59]:

\begin{exe}
    \ex \glll
    gilbil smex ga'a'y  \\
    gilbil smex ga'a-'y \\
    two bear see-\Fsg{}.\Ii{}  \\
    \glt "I saw two bears."  \hfill HH.\ElicitationDate{} (SF)
\end{exe}

[1:20:07]: The question here with \textit{=hl} is \textit{g̱a'ande} `where'...

\begin{exe}
    \ex \glll
    gilbilhl smex ga'a'y  \\
    gilbil=hl smex ga'a-'y \\
    two=\Cn{} bear see-\Fsg{}.\Ii{}  \\
    \glt "I saw two bears."  \hfill HH.\ElicitationDate{} (SF)
\end{exe}

[1:21:06]: Note that here HH says that \textit{g̱a'a} and \textit{sun} are two separate words. Idk if this is an intuition or if there is variable spelling for \textit{g̱a'asun}.

\begin{exe}
    \ex \glll
    gilbil(hl) smex ga'a'y g̱a'asun  \\
    gilbil(=hl) smex ga'a-'y g̱a'a[-t]-sun \\
    two(=\Cn{}) bear see-\Fsg{}.\Ii{} \Loc{}[-3.\Ii{}]=\Pn{}=\Demprox{} \\
    \glt "I saw two bears over here."  \hfill HH.\ElicitationDate{} (SF)
\end{exe}

[1:21:20]: Here it isn't finished without the additional location. Not really clear where this comes from.

\begin{exe}
    \ex \glll
    gilbilhl smexhl ga'a'y *(g̱a'asun) \\
    gilbil=hl smexhl ga'a-'y *(g̱a'a[-t]-sun) \\
    two=\Cn{} bear see-\Fsg{}.\Ii{} *(\Loc{}[-3.\Ii{}]=\Pn{}=\Demprox{} ) \\
    \glt "I saw two bears over here."  \hfill HH.\ElicitationDate{} (SF)
\end{exe}

[1:22:38]: with \textit{lax̱} is specifically `\textbf{on} the mountain'. Without it is ``up the road''.

\begin{exe}
    \ex \glll
    gilbilhl smex ga'a'y g̱a'ahl (lax̱) sg̱a'nist\\
    gilbil=hl smexhl ga'a-'y g̱a'a[-t]=hl (lax̱) sg̱a'nist \\
    two=\Cn{} bear see-\Fsg{}.\Ii{} \Loc{}[-3.\Ii{}]=\Cn{} (on) mountain \\
    \glt "I saw two bears up (on) the mountain."  \hfill HH.\ElicitationDate{} (VF)
\end{exe}

[1:24:13]: \textbf{Context:} Your are organizing a get-together, and you know that many people will come.

The stories glosses \textit{begwit} with \Tr{}, but arriving isn't actually transitive?? Are we sure it doesn't just sound like a transitive vowel because of the rounded \textit{gw} + \textit{-t} ?

\begin{exe}
    \ex \glll
    hilthl gat dim begwit \\
    hilt=hl gat dim begw-t \\
    many=\Cn{} person \Prosp{} come-\Tsg{}.\Ii{} \\
    \glt "A lot of people are going to come." \hfill HH.\ElicitationDate{} (VF)
\end{exe}

[1:24:55]: \textbf{Context:} [following above] now I don't know how many people will come to the party, but probably a lot.

At [1:26:18] I add the context of ``you are thinking to yourself'' \textit{ha'niig̱oodii'y}, so here Hector accepts it more readily.

\begin{exe}
    \ex \glll
    hildimaahl gat dim begwit \\
    hilt=imaa=hl gat dim begw-t \\
    many=\Epis{}=\Cn{} person \Prosp{} come-\Tsg{}.\Ii{} \\
    \glt "Maybe a lot of people are going to come." \hfill HH.\ElicitationDate{} (SF)
\end{exe}

[1:25:30]: In response to the above

\begin{exe}
    \ex
    \begin{xlist}
        \item \glll
        dimis g̱abii gat dim begwit \\
        dimis gabii gat dim begw-t\\
        dimis \Cnt{} perosn \Prosp{} come-\Tsg{}.\Ii{} \\
        \glt "How many people are gonna come?" \hfill HH.\ElicitationDate{} (VF)
        \item \glll
        hildimaa \\
        hilt=imaa \\
        many=\Epis{} \\
        \glt "Maybe a lot." \hfill HH.\ElicitationDate{} (VF)
    \end{xlist}
\end{exe}

[1:27:12]: 

\begin{exe}
    \ex \glll
    ha'niig̱oodi'y dim hilthl gat dim begwit \\
    ha-'nii-g̱oot-'y dim hilt=hl gat dim begw-t\\
    \Ins{}-on-heart-\Fsg{}.\Ii{} \Prosp{} many=\Cn{} person \Prosp{} come-\Tsg{}.\Ii{} \\
    \glt "I think a lot of people will come." \hfill HH.\ElicitationDate{} (VF)
\end{exe}

[1:27:45]: This is if you know for sure.

\begin{exe}
    \ex \glll
    ḵ'ap hilt dim begwit \\
    ḵ'ap hilt dim begw-t\\
    \Verum{} many \Prosp{} come-\Tsg{}.\Ii{} \\
    \glt "A lot of people will come." \hfill HH.\ElicitationDate{} (VF)
\end{exe}

[1:28:40]: I originally try with \textit{gilbilhl gat}, but Hector repeats it with the following. Interesting that the \textit{=hl} doesn't pattern identically to \textit{hilt} above.

This sentence doesn't actually mean living people, because of the inanimate numeral I used. 

\begin{exe}
    \ex \glll
    gilbil gathl dim begwit \\
    gilbil gat=hl dim begw-t \\
    two person=\Cn{} \Prosp{} come-\Tsg{}.\Ii{} \\
    \glt "Two people will come (=will be born)." \hfill HH.\ElicitationDate{} (VF)
\end{exe}

[1:30:30]:

\begin{exe}
    \ex \glll
    bag̱adilhl gat dim begwit \\
    bag̱adilhl gat dim begw-t \\
    two.\Hum{}=\Cn{} person= \Prosp{} come-\Tsg{}.\Ii{} \\
    \glt "Two people will come." \hfill HH.\ElicitationDate{} (SF)
\end{exe}

[1:31:03]: This is possible, but ``it'll be a question''. Repeated at [1:31:38], where it is just ``out there''... I assume it is still grammatical. 

At [1:33:50] there is reconfirmation that ``it works''.

Also at [1:35:08] I confirm that this doesn't have to be a question, it can be a statement.

\begin{exe}
    \ex \glll
    bag̱adilimaahl gat dim begwit \\
    bag̱adil=imaa=hl gat dim begw-t \\
    two.\Hum{}=\Epis{}=\Cn{} person \Prosp{} come-\Tsg{}.\Ii{} \\
    \glt "Maybe two people will come." \hfill HH.\ElicitationDate{} (SF)
\end{exe}

[1:31:25]: This is also in a conversational context.

\begin{exe}
    \ex \glll
    bag̱adilhl gat dim begwitimaa \\
    bag̱adil=hl gat dim begw-t=imaa \\
    two.\Hum{}=\Cn{} person \Prosp{} come-\Tsg{}.\Ii{}=\Epis{} \\
    \glt "Maybe two people will come." \hfill HH.\ElicitationDate{} (VF)
\end{exe}

[1:32:55]: Andrew asks in a context where you know two people are coming, how do you say ``maybe more people are coming''.

For \textit{wilaayi'y} I hear more of a \textit{h} or \textit{\emptyset-} onset.

HH translates as: ``I know there's two coming, but the other three I don't know''.

\begin{exe}
    \ex \glll
    wilaayi'y dim bekwhl bag̱adilt, ii nee diin wilax dim bekwhl hlag̱ats'uut \\
    wilaayi-'y dim bekw=hl bag̱adil-t, ii nee dii=n wilax dim bekwhl hla-g̱at-s'uu-t\\
    know-\Fsg{}.\Ii{} \Prosp{} come=\Cn{} two.\Hum{}-\Tsg{}.\Ii{}, \Ccnj{} \Neg{} \Foc{}=\Fsg{}.\I{} know \Prosp{} come=\Cn{} \Part{}-\Distr{}-other-\Tsg{}.\Ii{} \\
    \glt "I know that two (people) will come, but I don't know whether the others will come." \hfill HH.\ElicitationDate{} (VF)
\end{exe}

[1:33:10]:

\begin{exe}
    \ex \glll
    dim begwimaa 'nidiit \\
    dim bekw=imaa 'nidiit \\
    \Prosp{} come=\Epis{} \Tpl{}.\Iii{} \\
    \glt "They may come." \hfill HH.\ElicitationDate{} (VF)
\end{exe}